% Fuente: Auxiliar 1, P1.
Un estado desea planificar su capacidad eléctrica para los próximos $T$ años. El estado dispone de un pronóstico de $d_t$ MW, que se presume exacto, de la demanda de electricidad durante el año $t = 1, \dots, T$. La capacidad existente, representada por plantas alimentadas por petróleo que no se retirarán y estarán disponibles durante el año $t$, es $e_t$. Existen dos alternativas para ampliar la capacidad eléctrica: plantas de carbón o centrales nucleares. Existe un costo de capital de $c_t$ por MW de capacidad de carbón que entra en funcionamiento a principios del año $t$. El costo de capital correspondiente para las centrales nucleares es $n_t$. Por diversas razones políticas y de seguridad, se ha decidido que no más del 20\% de la capacidad total debe ser nuclear en ningún momento. Las plantas de carbón tienen una vida útil de 20 años, mientras que las nucleares duran 15 años. Se desea un plan de expansión de capacidad de costo mínimo.
